% ID: SEL-D0051-OB-E01
% ESTADO: cotejada. Segunda pasada visual y borrador compilado revisados.
% Fuente: Selectividad/sel_2013_matematicas_aplicadas/septiembre/reserva a/EXAMEN 3 -ANDALUCÍA.pdf
% SHA-256: 35f960392892cf4aa4b4a3e5b0a7b6150b323d584bb033f75db6a182b19206d4
% Página: 2; opción B; ejercicio original 1.

Se considera el recinto $R$ del plano determinado por las siguientes inecuaciones:
\[
5x-4y\leq20;\quad x+8y\leq48;\quad x\geq2;\quad y\geq0.
\]
\begin{enumerate}[label=\alph*)]
\item \textbf{(1.5 puntos)} Represente gráficamente el recinto $R$ y calcule sus vértices.
\item \textbf{(0.5 puntos)} Halle los valores máximo y mínimo que alcanza la función $F(x,y)=2x+12y$ en este recinto e indique dónde se alcanzan.
\item \textbf{(0.5 puntos)} Razone si existen valores $(x,y)$ pertenecientes al recinto para los que $F(x,y)=100$.
\end{enumerate}
